% Original arXiv 1611.02108v1, not the later published variant.
\section{The original CCHM cubical theory}
\label{sec:p07-cchm2016}

The theory considered here is Cohen, Coquand, Huber and M\"ortberg's
original November 2016 preprint, arXiv:1611.02108v1
\cite{CCHM2016ArxivV1}.  Its interval operations, face restrictions,
composition and Glue construction form one particular computational
presentation.  They are not ordinary identity-type rules with cubical
names attached.  The account separates the original rules from derived
constructions and from the metatheorems requiring a further argument.
The rule sheet is \ref{sheet:p07-cchm2016-cubical}; the local construction
proofs are in the corresponding annex.

\subsection{Two sorts of variables and the two lattices}
\label{sec:p07-cchm2016-dimensions}

Ordinary variables \(x,y\) have types.  Dimension names \(i,j,k\) range
over a fixed discrete infinite set and are declared \(i:\mathbb I\).
The expression \(\mathbb I\) is a dimension sort, not a type in the
ordinary dependent-product formation rule.  In particular, a name
abstraction \(\langle i\rangle t\) is not a lambda abstraction
\(\lambda i:\mathbb I.t\).

The dimension expressions are
\[
 r ::= 0_{\mathbb I}\mid1_{\mathbb I}\mid i\mid1-r
       \mid r\wedge s\mid r\vee s .
\]
They form the free de Morgan algebra on the declared names: a bounded
distributive lattice with an involution interchanging meet and join,
and interchanging zero and one.  Thus
\[
 1-(r\wedge s)=(1-r)\vee(1-s),\qquad
 1-(r\vee s)=(1-r)\wedge(1-s),\qquad1-(1-r)=r.
\]
Neither \(i\wedge(1-i)=0_{\mathbb I}\) nor
\(i\vee(1-i)=1_{\mathbb I}\) is imposed.  The source discusses a Kleene
alternative, but that alternative is not added to this presentation.

The distinct face lattice \(\mathbb F\) is generated by
\((i=0),(i=1)\), with
\((i=0)\wedge(i=1)=0_{\mathbb F}\), distributivity and lattice units.
It has no face negation constructor:
\[
 \varphi ::=0_{\mathbb F}\mid1_{\mathbb F}\mid(i=0)\mid(i=1)
                       \mid\varphi\wedge\psi\mid\varphi\vee\psi.
\]
In particular \((i=0)\vee(i=1)\ne1_{\mathbb F}\) in the unrestricted
context.  This is the boundary of an interval, not a cover of its
interior.  The notation \((r=1)\) is induced by the lattice map
\(i\mapsto(i=1)\), \(1-i\mapsto(i=0)\); \((r=0)\) means \((1-r=1)\).
Name substitution acts on both lattices through these maps.

A context can contain ordinary declarations, fresh name declarations
and face restrictions:
\[
 \Gamma ::=()\mid\Gamma,x:A\mid\Gamma,i:\mathbb I\mid\Gamma,\varphi .
\]
Judgments in \(\Gamma,\varphi\) impose \(\varphi=1_{\mathbb F}\).
For example, in \(\Gamma,(i=0)\), substitution of \(0\) for \(i\)
computes dimension expressions as expected.  A restriction is not a
hypothesis of a new propositional type and has no proof variable.
The source's five main judgments are context formation, type formation,
term typing, type equality and typed term equality.  Equalities are
judgmental, with equivalence and congruence closure; path inhabitation
does not imply judgmental equality.

We retain the source's substitution convention: \(t(i/r)\) replaces
the name \(i\) by \(r\), and \(B(x/a)\) replaces the ordinary variable
\(x\) by \(a\).  Substitution is simultaneous and capture-avoiding.
A typed substitution \(\Delta\vdash\sigma:\Gamma\) assigns an ordinary
term to each ordinary declaration and a dimension expression to each
name.  To substitute into \(\Gamma,\varphi\), it must additionally
satisfy \(\Delta\vdash\varphi\sigma=1_{\mathbb F}\).  Binder extension
\((\sigma,i/j)\) uses a name \(j\) fresh for the target context.

The core raw expressions comprise variables, dependent products and
sums, lambda abstraction and application, pairs and projections, the
natural-number type and its zero, successor and recursor, path types,
name abstraction and dimension application, systems, composition,
Glue with glue and unglue, and the universe.  Types and terms share
the raw-expression category in the Russell presentation; their
judgments distinguish their roles.  Expression application to an
ordinary term and application to a dimension expression have different
typing rules.  Bound ordinary variables and names are taken modulo
alpha conversion, with their distinct freshness requirements.

\begin{hbexample}{def:p07-cchm2016-face-substitution}{A connection substitution changes the extent}
The name substitution \(\sigma=(i/j\wedge k)\) from the context with
\(j,k:\mathbb I\) to the context with \(i:\mathbb I\) sends the face
\((i=0)\) to
\[
 ((j\wedge k)=0)=(j=0)\vee(k=0),
\]
and sends \((i=1)\) to
\[
 ((j\wedge k)=1)=(j=1)\wedge(k=1).
\]
The first equation uses the de Morgan involution before the interval
expression is mapped into the face lattice.  Thus a tube on the first
face becomes a tube on the union of two edges, while a tube on the
second becomes a tube on their common opposite corner.  Neither
becomes a total tube merely because opposite endpoints exist.
If the composition direction is a separate fresh name \(\ell\),
uniform substitution preserves the cap equality on precisely this
substituted extent; it cannot discard it or replace it by a Boolean
case split.  This exhibits why dimensions, faces and typed
substitutions must be recorded together.
\end{hbexample}

\subsection{Paths are cubes with specified ends}
\label{sec:p07-cchm2016-paths}

For \(A\) in \(\Gamma\), the type \(\mathsf{Path}\,A\,a\,b\) consists
of paths with endpoints \(a,b:A\).  Its original introduction and
elimination operations are name abstraction and application to a
dimension expression.  In the introduction rule \(A\) does not
depend on the newly bound name; the body \(t\) may.  The equations are
\[
 (\langle i\rangle t)\,r=t(i/r),\qquad
 p\,0=a,\quad p\,1=b,\qquad \langle i\rangle p\,i=p.
\]
The last is typed path extensionality.  These equations are not the
identity eliminator \(J\).

For a fixed \(A\), reflexivity is
\(1_a=\langle i\rangle a\), reversal is
\(\langle i\rangle p(1-i)\), and a connection square is
\(\langle i\rangle\langle j\rangle p(i\wedge j)\).
The four edges of its body are, respectively,
\[
 p(0\wedge j)=a,\quad p(1\wedge j)=p\,j,\quad
 p(i\wedge0)=a,\quad p(i\wedge1)=p\,i.
\]
For \(f:A\to B\), the action on a path is
\(\langle i\rangle f(p\,i)\).  For
\(p:(x:A)\to\mathsf{Path}\,B\,(f\,x)\,(g\,x)\),
function extensionality is
\(\langle i\rangle\lambda x:A.p\,x\,i\); its endpoints are \(f,g\)
by function eta.  These are explicit terms, not newly postulated axioms.

\begin{hblemma}{thm:p07-cchm2016-singleton}
 {proved=annex:p07.cchm2016.singleton}{A cubical contraction of a path singleton}
For \(a:A\), the type
\((x:A)\times\mathsf{Path}\,A\,a\,x\) has center \((a,1_a)\).
Its contraction at \((b,p)\) is
\[
 \langle i\rangle
       (p\,i,\langle j\rangle p(i\wedge j)).
\]
\end{hblemma}

\subsection{Systems, extent and composition}
\label{sec:p07-cchm2016-composition}

A system \([\varphi_1\,t_1,\ldots,\varphi_n\,t_n]\) can glue
locally defined terms only when its faces cover the current context
and its branches agree on every overlap.  A tube of extent \(\varphi\)
is a term in the restricted context \(\Gamma,\varphi,i:\mathbb I\);
it need not extend to \(\Gamma,i:\mathbb I\).  The abbreviation
\(a:A[\varphi\mapsto u]\) states \(a:A\) and the equality \(a=u:A\)
under \(\varphi\).  It is not a new type former.

Given
\[
 \Gamma,i:\mathbb I\vdash A,\quad
 \Gamma,\varphi,i:\mathbb I\vdash u:A,\quad
 \Gamma\vdash a_0:A(i/0)[\varphi\mapsto u(i/0)],
\]
with \(\varphi\) in \(\Gamma\), composition produces
\[
 \mathsf{comp}^{i}A[\varphi\mapsto u]a_0
       :A(i/1)[\varphi\mapsto u(i/1)].
\]
Thus the side is independent of the bound composition direction, and
the bottom agrees with it on that side.  On full extent composition
returns the top of the tube.  Uniformity means that this operation
commutes with substitutions, after capture-avoiding renaming of its
bound dimension.  Empty brackets abbreviate extent \(0_{\mathbb F}\),
not an omitted compatibility condition on a nonempty tube.

Filling is derived rather than added as another primitive:
\[
 \mathsf{fill}^{i}A[\varphi\mapsto u]a_0
 =\mathsf{comp}^{j}A(i/i\wedge j)
       [\varphi\mapsto u(i/i\wedge j),(i=0)\mapsto a_0]a_0,
\]
where \(j\) is fresh for \(\Gamma,i\).  This expression lives along
\(i\) in the original family \(A\).

\begin{hblemma}{thm:p07-cchm2016-filling}
 {proved=annex:p07.cchm2016.filling}{The three filling boundaries}
The displayed filling has bottom \(a_0\), top
\(\mathsf{comp}^{i}A[\varphi\mapsto u]a_0\), and restriction \(u\)
on \(\varphi\).  All three equalities are judgmental.
\end{hblemma}

A complete two-path composition, for \(p:a\leadsto b\) and
\(q:b\leadsto c\) in a fixed \(A\), is
\[
 p\mathbin{\cdot}q
 =\langle i\rangle\mathsf{comp}^{j}A
       [(i=0)\mapsto a,(i=1)\mapsto q\,j](p\,i).
\]
At the bottom the two side equations are \(p\,0=a\) and
\(p\,1=b=q\,0\).  The opposite \(i\)-faces have empty overlap.
At \(i=0\), full-face composition returns \(a\); at \(i=1\) it
returns \(q\,1=c\).  This proves the typing and endpoints without
assuming a groupoid law.

Transport is exactly
\(\mathsf{transp}^{i}A\,a_0=\mathsf{comp}^{i}A[]a_0\).
There is no additional general judgmental equation
\(\mathsf{transp}^{i}A\,a_0=a_0\) for a constant family in the
original rules.  Filling nevertheless gives a path from \(a_0\) to
that transport.  Judgmental and propositional regularity must therefore
be kept separate.

\subsection{The type-specific composition clauses}
\label{sec:p07-cchm2016-type-composition}

The computation of composition is determined at each type constructor.
For \(C=(x:A)\to B\), a cap \(h_0\), tube \(h\), and
\(z:A(i/1)\), first form the backward filling
\[
 w=\mathsf{fill}^{i}A(i/1-i)[]z,\qquad v=w(i/1-i).
\]
Then \(v:A\) along \(i\), with \(v(i/1)=z\), and the result applied
to \(z\) is
\[
 \mathsf{comp}^{i}B(x/v)[\varphi\mapsto h\,v]
                          (h_0\,v(i/0)).
\]
The reversal is essential: an arbitrary argument at the top is not
automatically an argument in the varying domain at the bottom.

For \(C=(x:A)\times B\), cap \(z_0\) and tube \(z\), set
\[
 a=\mathsf{fill}^{i}A[\varphi\mapsto z.1]z_0.1,\qquad
 c_1=a(i/1),\qquad
 c_2=\mathsf{comp}^{i}B(x/a)[\varphi\mapsto z.2]z_0.2.
\]
The composed pair is \((c_1,c_2)\).  The first component is filled,
not just composed, so that the second component has a family throughout
the composition direction.

For natural numbers, composition on a tube and cap both headed by
zero gives zero; on a tube \(\mathsf{s}\,u\) with cap
\(\mathsf{s}\,a_0\), it gives
\(\mathsf{s}(\mathsf{comp}^{i}\mathbb N
[\varphi\mapsto u]a_0)\).  These constructor equations alone are not
a proof that every closed natural-number term reaches a constructor.

For \(C=\mathsf{Path}\,A\,a\,b\), with varying \(A,a,b\), the clause is
\[
 \mathsf{comp}^{i}C[\varphi\mapsto p]p_0
 =\langle j\rangle\mathsf{comp}^{i}A
       [\varphi\mapsto p\,j,(j=0)\mapsto a,(j=1)\mapsto b](p_0\,j).
\]
The \(j\)-faces enforce the endpoints of the resulting path, and the
path tube's endpoint equations supply the overlaps with \(\varphi\).

\subsection{Contractibility and the lifting operation for equivalences}
\label{sec:p07-cchm2016-equivalence}

The source uses
\[
 \mathsf{isContr}(A)
 =(c:A)\times((x:A)\to\mathsf{Path}\,A\,c\,x),
\]
\[
 \mathsf{isEquiv}(T,A,f)
 =(y:A)\to\mathsf{isContr}((x:T)\times
                                  \mathsf{Path}\,A\,y\,(f\,x)),
\quad
 \mathsf{Equiv}(T,A)=(f:T\to A)\times\mathsf{isEquiv}(T,A,f).
\]
Notice the fiber path orientation \(y\leadsto f\,x\).
We write \(f.1\) for the map and \(f.2\) for its contraction data.

\begin{hblemma}{thm:p07-cchm2016-extension}
 {proved=annex:p07.cchm2016.extension}{Contractibility and uniform partial extension}
From \(c:\mathsf{isContr}(A)\), every partial \(u:A\) of extent
\(\varphi\) has the specified extension
\[
 \mathsf{contr}\,c[\varphi\mapsto u]
 =\mathsf{comp}^{i}A[\varphi\mapsto c.2\,u\,i]c.1.
\]
Conversely, an operation extending every such partial element in every
context extension gives contraction data for \(A\).
\end{hblemma}

Consequently, \(f:\mathsf{Equiv}(T,A)\), \(a:A\), and partial
\((t,p):(x:T)\times\mathsf{Path}\,A\,a\,(f.1\,x)\) have an extension
\[
 \mathsf{equiv}\,f[\varphi\mapsto(t,p)]a
 =\mathsf{contr}(f.2\,a)[\varphi\mapsto(t,p)].
\]
Both the point and its path extend strictly on the face.  This is the
lifting operation used by Glue, rather than an unexplained selection
of an inverse point.

For a varying map \(f:T\to A\), a tube \(t\) and compatible cap \(t_0\),
write \(v=\mathsf{fill}^{i}T[\varphi\mapsto t]t_0\).  The source's
composition-preservation path is
\[
 \mathsf{pres}^{i}f[\varphi\mapsto t]t_0
 =\langle j\rangle\mathsf{comp}^{i}A
       [\varphi\vee(j=1)\mapsto f\,v](f(i/0)\,t_0).
\]
It goes from
\(\mathsf{comp}^{i}A[\varphi\mapsto f\,t](f(i/0)t_0)\)
to \(f(i/1)(\mathsf{comp}^{i}T[\varphi\mapsto t]t_0)\).
On \(\varphi\) it is constant.
\begin{hblemma}{thm:p07-cchm2016-preservation}
 {proved=annex:p07.cchm2016.preservation}{Typing and faces of the preservation path}
The displayed term has precisely those endpoints and the stated
constant restriction on the tube face.
\end{hblemma}

\subsection{Glue and its actual composition algorithm}
\label{sec:p07-cchm2016-glue}

For \(A\) in \(\Gamma\), and \(T\) and \(f:\mathsf{Equiv}(T,A)\)
defined under \(\varphi\), form
\[
 G=\mathsf{Glue}[\varphi\mapsto(T,f)]A.
\]
Under \(\varphi\), \(G=T\) judgmentally.  A point is introduced as
\(\mathsf{glue}[\varphi\mapsto t]a\) with
\(a:A[\varphi\mapsto f.1\,t]\).  The map \(\mathsf{unglue}:G\to A\)
has restriction \(f.1\), and satisfies the beta and eta equations in
the sheet.  No equation \(\mathsf{Glue}[]A=A\) is included.  An
equivalence to \(A\) is not a judgmental collapse to \(A\).

Now allow the gluing face itself to depend on the composition
direction:
\(G(i)=\mathsf{Glue}[\varphi(i)\mapsto(T(i),f(i))]A(i)\).
For a tube \(b\) of extent \(\psi\), independent of \(i\), and cap
\(b_0\), set \(a=\mathsf{unglue}\,b\), \(a_0=\mathsf{unglue}\,b_0\).
The operation \(\forall i\) on faces sends both \(i\)-face generators
to zero and fixes the other generators, extending as a lattice
homomorphism.  Put \(\delta=\forall i\,\varphi\); then
\(\delta\leq\varphi\) throughout the direction.  Define
\[
 a'_1=\mathsf{comp}^{i}A[\psi\mapsto a]a_0,
 \qquad
 t'_1=\mathsf{comp}^{i}T[\psi\mapsto b]b_0
                     \quad\hbox{under }\delta ,
\]
and, also under \(\delta\),
\(\omega=\mathsf{pres}^{i}(f.1)[\psi\mapsto b]b_0\).
This path goes from \(a'_1\) to \(f(i/1).1\,t'_1\).
Under \(\varphi(i/1)\) use the fiber extension
\[
 (t_1,\alpha)=
 \mathsf{equiv}\,f(i/1)
 [\delta\mapsto(t'_1,\omega),
  \psi\mapsto(b(i/1),\langle j\rangle a'_1)]\,a'_1 .
\]
Finally set
\[
 a_1=\mathsf{comp}^{j}A(i/1)
       [\varphi(i/1)\mapsto\alpha\,j,\psi\mapsto a(i/1)]a'_1,
 \qquad
 b_1=\mathsf{glue}[\varphi(i/1)\mapsto t_1]a_1 .
\]
This is the original Glue composition, including the correction from
\(a'_1\) to \(a_1\).  Merely composing in \(A\) and choosing an
arbitrary preimage would miss the required boundary equations.

\begin{hblemma}{thm:p07-cchm2016-face-decomposition}
 {proved=annex:p07.cchm2016.face-decomposition}{The face decomposition used by Glue}
For every face expression,
\[
 \varphi=(\forall i\,\varphi)
       \vee(\varphi\wedge(i=0))\vee(\varphi\wedge(i=1)).
\]
Moreover \(\forall i\,\varphi\) is the largest \(i\)-independent face
below \(\varphi\), and
\(\varphi\wedge(i=\varepsilon)\leq\varphi(i/\varepsilon)\).
\end{hblemma}

For the source's explicit face
\[
 \varphi=(i=0)\vee((i=1)\wedge(j=0))\vee(j=1),
\]
the operation gives \(\forall i\,\varphi=(j=1)\).
It deletes the \(i=0\) summand and the summand containing \(i=1\),
but keeps the independent \(j=1\) summand.  In the Glue algorithm,
this means that the partial type \(T\) can be composed throughout
the whole \(i\)-direction on \(j=1\).  Its presence only at one
\(i\)-endpoint would not suffice for the intermediate \(T\)-filling.

\begin{hblemma}{thm:p07-cchm2016-glue-composition}
 {proved=annex:p07.cchm2016.glue-composition}{Typing and boundary equations of Glue composition}
The displayed \(b_1\) is well typed in \(G(i/1)\), equals \(b(i/1)\)
under \(\psi\), and equals \(t'_1\) under \(\delta\).  If the gluing
face is full, it reduces to the composition in \(T\).
\end{hblemma}

\begin{hbtheorem}{thm:p07-cchm2016-unglue}
 {proved=annex:p07.cchm2016.unglue}{Unglue has contractible fibers}
For the displayed \(G\), the map
\(\mathsf{unglue}:G\to A\) has a term of
\(\mathsf{isEquiv}(G,A,\mathsf{unglue})\), constructed using
composition and fiber extension.
\end{hbtheorem}

\subsection{The universe: line equivalence before universe composition}
\label{sec:p07-cchm2016-universe}

The source adds a Russell universe \(U\): \(U\) is a type, and
\(A:U\) may be used as a type.  The ordinary type formers, paths,
natural numbers and Glue have versions producing elements of \(U\)
when their input types lie in \(U\).  It does not assume \(U:U\).

For a type line \(i:\mathbb I\vdash E\), let \(A=E(i/0)\),
\(B=E(i/1)\).  Forward and backward transport are
\[
 f=\lambda x:A.\mathsf{transp}^{i}E\,x,\qquad
 g=\lambda y:B.\mathsf{transp}^{i}E(i/1-i)\,y.
\]
They need not be judgmental inverses.  The filled transports are
\[
 u=\lambda x:A.\mathsf{fill}^{i}E[]x,\qquad
 v=\lambda y:B.(\mathsf{fill}^{i}E(i/1-i)[]y)(i/1-i).
\]
For these displayed families, \(u(0)\) and \(u(1)\) abbreviate
the substitutions \(u(i/0)\) and \(u(i/1)\), and likewise for \(v\).
They are not dimension applications of a term of function type.
Thus \(u(0)x=x,\ u(1)x=f\,x,\ v(0)y=g\,y,\ v(1)y=y\).
The original fiber contraction, given fully in the annex, proves
\(\mathsf{isEquiv}(A,B,f)\).  The resulting pair is denoted
\(\mathsf{equiv}^{i}E:\mathsf{Equiv}(A,B)\).

\begin{hbtheorem}{thm:p07-cchm2016-line-equivalence}
 {proved=annex:p07.cchm2016.line-equivalence}{A type line yields an equivalence}
The forward transport \(f\) has contractible fibers with the source's
path orientation \(y\leadsto f\,x\).  Its equivalence structure is
constructed, not assumed.
\end{hbtheorem}

For \(E:U\) under \(\varphi,i\) and \(A_0:U\) agreeing with \(E(0)\)
under \(\varphi\), universe composition is
\[
 \mathsf{comp}^{i}U[\varphi\mapsto E]A_0
 =\mathsf{Glue}
 [\varphi\mapsto(E(i/1),\mathsf{equiv}^{i}E(i/1-i))]A_0 .
\]
The reversed line supplies an equivalence from \(E(1)\) to
\(E(0)=A_0\), which is the direction required by Glue formation.
On \(\varphi\), the full-face Glue rule returns \(E(1)\);
outside it the constructed type need not judgmentally equal \(A_0\).
This checks the universe composition's type and its boundary; it is
not yet a proof of univalence.

An equivalence \(e:\mathsf{Equiv}(A,B)\) also gives a type path
\[
 \mathsf{eqToPath}(e)=
 \langle i\rangle\mathsf{Glue}
 [(i=0)\mapsto(A,e),(i=1)\mapsto(B,\mathsf{equiv}^{k}B)]B .
\]
Here the second equivalence is the one constructed from the constant
type line, as in source Appendix B; it is not silently replaced by a
judgmentally identical identity equivalence.  The opposite endpoint
faces are disjoint, and the Glue face equations give the ends \(A,B\).
At \(i=0\) and \(i=1\) this expression is thus a genuine type path,
not an axiom asserting that every equivalence already is equality.

\begin{hbtheorem}{thm:p07-cchm2016-univalence}
 {cited=CCHM2016ArxivV1}{Source univalence construction}
The source proves univalence for its universe using Glue and the
constructed equivalence operations.  In particular, the canonical
map from type paths to equivalences is itself an equivalence.
\end{hbtheorem}

\begin{proofboundary}{def:p07-cchm2016-univalence-boundary}
 {What remains between the type path and full univalence}
The local proofs above construct line equivalence, Glue composition,
and the equivalence of unglue.  They do not by themselves prove that
the canonical path-to-equivalence map is an equivalence.  Source
Section 7.2 and Appendix B further require the transport comparison
\(\mathsf{transp}(\mathsf{eqToPath}(e))\leadsto e.1\), the
two-dimensional Glue coherence comparing a type path with the
reconstruction of its line equivalence, and the passage to full
equivalence data, not just function parts.  The source uses the
propositionality of \(\mathsf{isEquiv}\) in one route and refers to a
formal development for details in another.  Those full comparison
and coherence arguments are not supplied here.  The cited theorem's
standing must not be promoted to a locally proved theorem by treating
the missing coherence as an axiom.
\end{proofboundary}

\showcalculus{p07-cchm2016-cubical}

\subsection{Presheaf meaning and the computational theorem boundary}
\label{sec:p07-cchm2016-semantics}

In the source model a cube is a finite set of names \(I\), and a map
\(J\to I\) assigns to each name in \(I\) an expression in the free
de Morgan algebra on \(J\).  Cubical sets are contravariant functors
on this category.  A semantic type over a cubical context \(\Gamma\)
has fibers \(A\rho\) and functorial restrictions; a term is a natural
section.  Type substitution pulls back these fibers and operations.
The face restriction \(\Gamma,\varphi\) contains exactly the
environments \(\rho\) at which \(\varphi\rho=1_{\mathbb F}\).

The critical additional datum is a \emph{uniform} composition
structure.  Given a direction \(i\notin I\), an environment over
\(I,i\), a face in \(\mathbb F(I)\), a natural partial tube and a
compatible cap, it supplies a top point.  For every cube map
\(h:J\to I\), restriction of that point must equal composition of
the restricted tube and cap, with a fresh renamed direction.
Full-face composition must return the tube's top.  Existence of one
filler for each isolated box is not enough; the restriction equations
are part of the structure.

Semantic Glue retains \(T\rho\) itself where the face is full, and
otherwise stores a partial \(T\)-section together with an \(A\)-point
whose image agrees on the face.  Its restrictions may move from the
stored-pair case to the full-face \(T\)-case.  Equivalence structure
supplies a uniform extension of partial fiber data, permitting the
same corrected composition algorithm.  The semantic universe consists
of small fibrant types over representable cubes.  Interpreting the
Russell syntax uses two ambient Grothendieck universes, described in
the source as tiny and small; this is an external set-theoretic model
assumption, not the inconsistent judgment \(U:U\).

\begin{hbtheorem}{thm:p07-cchm2016-soundness}
 {cited=CCHM2016ArxivV1}{The source's interpretation theorem}
The source's Section 8 interpretation sends derivable contexts to
cubical sets, derivable types to semantic universe elements, terms
to natural sections, substitutions to natural transformations, and
judgmental equalities to semantic equalities.
\end{hbtheorem}

That statement does not establish normalization or decidability of
the syntactic conversion problem.  The original paper's concluding
discussion treats a normalization proof and the resulting decidability
argument as further work.  Its citation for canonicity is instead
Huber's separate July 2016 preprint, arXiv:1607.04156v1
\cite{Huber2016CanonicityV1}.  The separately stated canonicity scope is
natural-number terms in contexts containing only dimension names:
they are judgmentally equal to numerals.  This includes the closed
case, but does not say that an arbitrary \(x:\mathbb N\) is a numeral.
The separate paper's typed operational semantics and presheaf-like
computability proof are not reconstructed by the present rule sheet.

\begin{proofboundary}{def:p07-cchm2016-metatheory-boundary}
 {Semantic soundness, canonicity and normalization remain distinct}
The full derivation induction for semantic soundness, including all
substitution and naturality cases and equality of the associated
composition structures, is source-cited rather than locally proved.
The separate Huber canonicity proof has not been read in full or
reconstructed here.  No normalization, decision procedure or full
canonicity proof is claimed.  The source's circle, truncation and
identity-type extensions are also not reconstructed in this account.
\end{proofboundary}

Compared with an axiomatic homotopy type theory, the important
difference is not that both speak of paths and equivalences.  Here
dimensions and face formulas control a concrete composition operation;
line equivalence and type-path construction use that operation and
Glue.  Their boundaries compute by specified equations.  Axiomatic
univalence, by contrast, does not supply these original cubical
grammar, face and composition clauses merely by being assumed.
Conversely, the clauses do not license importing every identity-type
computation rule, a normalization theorem, or a later cubical
presentation without its own construction and proof.
